(a)\ 求大灰熊的身体观测数据的六个总体均值的大样本 $95\%$ 联合置信区间.

(b)\ 求体重和腰围的均值的大样本 $95\%$ 联合置信椭圆.

(c)\ 求在 (a) 中的六个均值的大样本 $95\%$ 庞弗罗尼置信区间.

(d)\ 参考 (b), 对 $m=6$ 构造重量和腰围均值的 $95\%$ 庞弗罗尼置信矩形, 并和 (b) 中的置信椭圆进行比较.

(e)\ 对 $m=6+1=7$(所有单个均值的表述), 求头宽均值—头长均值的 $95\%$ 庞弗罗尼置信区间.

\noindent 5.10\quad 参考例 1.10 (见\hyperref[tab:1.4]{表 1.4}) 中的熊的成长数据. 关注体长观测值.

(a)\ 求四个总体的体长均值的 $95\%$ $T^2$ 联合置信区间.

(b)\ 参考 (a), 求体长均值的相继三年的年增长的 $95\%$ $T^2$ 联合置信区间.

(c)\ 求从 2 岁到 3 岁的体长的增长均值以及 4 岁到 5 岁的体长的增长均值的 $95\%$ $T^2$ 置信椭圆.

(d)\ 参考 (a) 和 (b), 构造包含四个体长均值以及体长均值的相继三年的年增长的集合的 $95\%$ 庞弗罗尼置信区间.

(e)\ 参考 (c) 和 (d), 将从 2 岁到 3 岁的体长的增长均值以及 4 岁到 5 岁的体长的增长均值的 $95\%$ 庞弗罗尼置信矩形和与 $T^2$ 方法得到的置信椭圆进行比较.

\noindent 5.11\quad 一名人类学家对 9 根古秘鲁人头发的头发完成了矿物分析, 测得其中的铬(Cr)含量 $x_1$ 和锶(Sr)含量 $x_2$(单位为 p.p.m, 即百万分之一)分别为

\begin{center}
\begin{tabular}{lccccccccc}
\toprule
$x_1$(Cr) & 0.48 & 40.53 & 2.19 & 0.55 & 0.74 & 0.66 & 0.93 & 0.37 & 0.22\\
$x_2$(Sr) & 12.57 & 73.68 & 11.13 & 20.03 & 20.29 & 0.78 & 4.64 & 0.43 & 1.08\\
\bottomrule
\end{tabular}
\end{center}

已知铬含量较低(低于或等于 0.100 p.p.m)可能是糖尿病的表现, 而锶含量的高低则说明动物蛋白摄入的多少.

(a)\ 假定这 9 根秘鲁人头发代表一个来自属于某种特定古秘鲁文化的随机样本. 试构造并画出总体均值向量 $\boldsymbol{\mu}'=[\mu_1,\mu_2]$ 的 $90\%$ 联合置信椭圆.

(b)\ 用 (a) 中构造的椭圆向每一坐标轴投影, 求出 $\mu_1$ 与 $\mu_2$ 的 $90\%$ 联合置信区间(可利用结论 5.3). 是否有迹象表明这种古秘鲁人头发的平均锶含量为 10? 也就是说, 所有的点(任意 $\mu$ 值, 10)是否都在置信域中? $[0.30,10]$ 是否为 $\boldsymbol{\mu}$ 的真值? 试作讨论.

(c)\ 这些数据是否显示出二元正态性? 利用 Q-Q 图与散布图来考察这个问题. 若数据没有二元正态性, 那么 (a) 与 (b) 中的结论又意味着什么?

(d)\ 摒弃明显的离群观测值后重新进行分析, 推断是否会有所改变? 试作评论.

\noindent 5.12\quad 设有若干个分量缺损的观测数据为
\[
\mathbf{X}=\begin{bmatrix}
3&6&0\\
4&4&3\\
-&8&3\\
5&-&-
\end{bmatrix}
\]
试用 5.7 节中的预测-估计算法估计 $\boldsymbol{\mu}$ 和 $\boldsymbol{\Sigma}$, 要求给出其初始估计值及一步迭代的修正估计值.

\noindent 5.13\quad 对\hyperref[tab:5.1]{表 5.1} 的汗液数据确定 $-n\ln(|\hat{\boldsymbol{\Sigma}}|/|\hat{\boldsymbol{\Sigma}}_0|)$ 的近似分布(参见结论 5.2).

\noindent 5.14\quad 试建立一个类似于\hyperref[tab:5.4]{表 5.4} 的新表, 表中的元素为(单一 $t$ 区间)的宽度/(庞弗罗尼 $t$
